\chapter{Introducción}
\label{chapter:1-Introducción} 

%\textit{Aquí se motiva el trabajo, se plantea y define el problema, se deja claro cuales son los objetivos (general del proyecto, si correspondiese o si está inmerso en un proyecto de mayor alcance, y los específicos), se plantean los resultados esperados, se establecen resumidamente las conclusiones y se describe la organización general del documento.
% Este y los siguientes capítulos pueden incluir referencias, algunos ejemplos de referencias pueden ser \cite{CitekeyArticle}, \cite{CitekeyBook}, \cite{CitekeyInproceedings}, \cite{CitekeyManual} y \cite{CitekeyMisc}. Se sugiere citar usando el formato APA.}
%En la introducción se deberá:
%Explayar el resumen previo. 
%Explicar en detalle algunas de las %posibles falencias que se basan en %utilizar modelos que no fueron entrenados %con contenido educativo o que no fueron %entrenados con contenido local/regional.
%Ej.:
%nombres, enfoques o convenciones distintos %que nosotros (UdelaR o Uruguay) podamos %usar en comparación con un ámbito %internacional.
%Un caso de valor puede ser que el Método %Simplex (Introducción a la Investigación %Operativa) se dé de manera distinta.
%
Se deberá detallar el scope del proyecto, el alcance que nos propusimos.
Se deberá detallar cada uno de los parámetros de entrada y de los parámetros que se investigaron y evaluaron para llegar a un mejor resultado.
Se deberá detallar en base a qué evaluamos un modelo / un resultado como mejor o peor que otro.

--Lo Siguiente ya estaba escrito--

El acceso a información confiable y relevante es un desafío constante para los estudiantes universitarios. Aunque las grabaciones de clase pueden ser una herramienta útil, presentan varias limitaciones: su consumo es lento, pueden contener problemas de audio que dificultan la comprensión, y en muchos casos, solo una fracción del contenido resulta realmente valiosa. Además, la información presentada en clase puede quedar desactualizada con el tiempo, lo que genera confusión y dificulta el aprendizaje efectivo.

Este trabajo de tesis se centra en el desarrollo de un producto basado en el fine-tuning de un modelo de lenguaje (LLM) con contenido educativo específico de una facultad, utilizando como fuente principal las grabaciones de las clases. El objetivo es ofrecer a los estudiantes una herramienta en formato de pregunta-respuesta que les permita acceder rápidamente a información precisa y relevante, directamente basada en sus propias clases.

La implementación de esta solución no solo agiliza el proceso de aprendizaje, sino que también proporciona un medio más eficiente para digerir la información. Los estudiantes podrán realizar consultas en tiempo real sobre el contenido de sus cursos, sin necesidad de revisar extensos videos o depender de fuentes externas cuya fiabilidad puede ser incierta. De esta manera, se garantiza que el conocimiento adquirido provenga de una fuente validada y alineada con el plan de estudios de la facultad.

Este enfoque representa un avance significativo en la optimización del aprendizaje, brindando a los estudiantes una alternativa moderna y efectiva para complementar su formación académica.

\section{Objetivos del Proyecto}
\label{section:1-Objetivos del Proyecto}


\section{Resultados Obtenidos}
\label{section:1-Resultados Obtenidos}


\section{Organización del Documento}
\label{section:1-Organización del Documento}

El proximo documento se estructura de la siguiente manera:

\autoref{chapter:2-Marco Teórico} 
En este capítulo se presenta una revisión del estado del arte, los conceptos teóricos relevantes que sustentan el desarrollo del proyecto y una descripción para familiarizar al lector con OpenFing.

\autoref{chapter:3-Requisitos del Sistema} 
En este capítulo se formalizan los requisitos funcionales y no funcionales del sistema propuesto. 

\autoref{chapter:4-Arquitura y Diseño} 
En este capítulo se expone la arquitectura general del sistema y el diseño detallado de sus componentes. Se presentan diagramas de arquitectura y flujos de datos.

\autoref{chapter:5-Implementación} 
En este capítulo se describe las  implementación del sistema, las tecnologías utilizadas y sus expectativas.

\autoref{chapter:6-Experimentación} 
En este capítulo se presentan los experimentos realizados para evaluar el desempeño del sistema. Se describen los conjuntos de datos utilizados, las métricas de evaluación, y los resultados obtenidos. 

\autoref{chapter:7-Trabajo a futuro} 
En este capítulo se exploran las posibles direcciones futuras mejoras del sistema. 

\autoref{chapter:8-Conclusiones} 
En este capítulo se resumen los principales hallazgos del proyecto y se reflexiona sobre el cumplimiento de los objetivos planteados. 